% S5: ban TIENG VIET cua paper2/blocks/09_cost.tex.
% Chi duoc \input boi paper2_hujos/supplementary_vi.tex.
% \label{sec:cost} GIU LAI vi ban EN van \ref no.
% Chu trong \underbrace dich theo dung ten cot cua Bang S5 (tab_baseline):
% phan cung / ket noi / nhan cong / tham dinh.
\section*{S5. Mô hình chi phí và điểm hòa vốn}
\label{sec:cost}

Giả sử một hợp tác xã phục vụ $A$ héc-ta với $N$ vụ mùa trong đời thiết bị, và chi phí
MRV mỗi héc-ta mỗi vụ phân tích thành
\begin{equation}
C=\underbrace{\frac{n_{\mathrm{dev}}\,p_{\mathrm{dev}}}{A\,N}}_{\text{phần cứng}}
+\underbrace{c_{\mathrm{conn}}}_{\text{kết nối}}
+\underbrace{h\,w}_{\text{nhân công}}
+\underbrace{\frac{c_{\mathrm{audit}}}{A}}_{\text{thẩm định}},
\label{eq:cost}
\end{equation}
trong đó $n_{\mathrm{dev}}$ và $p_{\mathrm{dev}}$ là số lượng và đơn giá thiết bị, $h$
là số giờ công mỗi héc-ta mỗi vụ và $w$ là tiền công nông thôn. Doanh thu mỗi héc-ta là
\begin{equation}
R=\kappa\,\rho ,
\label{eq:revenue}
\end{equation}
với $\kappa$ là số tín chỉ cấp cho mỗi héc-ta và $\rho$ là giá mỗi tín chỉ. Phương án
khả thi về kinh tế khi
\begin{equation}
C\ \le\ \beta R ,\qquad \beta\in(0,1),
\label{eq:breakeven}
\end{equation}
và $\beta=0{,}5$ là mức mà giám sát ngốn hết một nửa doanh thu tín chỉ. Thế
\eqref{eq:revenue} vào ta được lượng tín chỉ tối thiểu để hòa vốn
\begin{equation}
\kappa_{\min}=\frac{C}{\beta\rho}.
\label{eq:kappamin}
\end{equation}

Các tầng khác nhau ở $n_{\mathrm{dev}}$ và, quan trọng hơn, ở chỗ chúng quan sát được
đại lượng nào trong ba đại lượng làm giả được của \eqref{eq:three}. Tầng~0 là tầng vật
lý: $n_{\mathrm{dev}}=0$ ngoài bộ điều khiển, nhật ký được ghi trên điện thoại mà nông
dân đã có sẵn, và các bộ lọc ở Mục~\ref{sec:filters} chạy trên dữ liệu thời tiết công
khai với chi phí biên bằng không. Tầng~1 thêm một đồng hồ lưu lượng tại trạm bơm của hợp
tác xã --- $n_{\mathrm{dev}}=1$ cho toàn bộ diện tích $A$ chứ không phải mỗi thửa một
cái --- và quan sát được tổng lượng nước. Tầng~2 thêm một nút đo mực nước cho mỗi hợp
tác xã, neo lại bản sao số và khôi phục tính hiệu chuẩn mà một hệ MRV dựa trên đo đạc
dù sao cũng cần.

Không tầng nào trong ba tầng đó quan sát được dấu thời điểm, và theo
Mệnh đề~\ref{prop:blind} thì hai tầng đầu bất biến chính xác đối với phép biến đổi sinh
lợi. Vì vậy neo thiết bị ở Mục~\ref{sec:anchor} không phải là một tầng thứ tư tuỳ chọn
mà là điều kiện để bằng chứng được chấp nhận: nó do chính bộ điều khiển đóng máy bơm
sinh ra, nên chỉ thêm việc lưu trữ và truyền đi một bản ghi vốn đã tồn tại, và không
thêm $n_{\mathrm{dev}}$ nào.

Bảng~\ref{tab:baseline} so sánh phương án này với ba kiến trúc mà nó thay thế. Với
$\kappa\in[4,6]$ và $\rho\in[\$15,\$25]$~\cite{rmi2025rice,upstream2026price}, doanh thu
là $R\in[\$60,\$150]$ mỗi héc-ta mỗi vụ, nên $\kappa_{\min}$ tại $\beta=0{,}5$ tương ứng
với một ngân sách giám sát $30$--$75$~USD mỗi héc-ta. Toàn bộ ngăn xếp của bài tốn
$8{,}4$ USD mỗi héc-ta mỗi vụ so với ngân sách $30$--$75$ đó, tức thấp hơn trần từ
$3{,}6$ đến $8{,}9$ lần; vì vậy biên độ được phát biểu thành một dải chứ không phải một
bậc độ lớn. Nó cũng không phải dòng rẻ nhất của Bảng~\ref{tab:baseline}: viễn thám kèm
kiểm tra thực địa tốn $8{,}0$, tức rẻ hơn $0{,}4$. Hai phương án trả lời hai câu hỏi
khác nhau --- phương án rẻ hơn xác nhận rằng một pha khô đã xảy ra nhưng không giữ bản
ghi theo từng lệnh nào để đối chiếu với một nhật ký bị làm giả --- và chi phí tăng thêm
của neo so với một bộ điều khiển đã lắp sẵn là khoảng $1\%$ doanh thu ở mức thận trọng.
Bảng~\ref{tab:baseline} phân tích phép so sánh với ba kiến trúc mà nó thay thế; hạng mục
kết nối được tách riêng vì đó chính là thứ phân biệt một thiết kế mỗi thửa một cảm biến
với một gateway dùng chung.

\begin{remark}
Giá thiết bị trong Bảng~\ref{tab:baseline} là ước tính phục vụ lập kế hoạch, không phải
báo giá. Thứ tự giữa các tầng và bất đẳng thức hòa vốn \eqref{eq:breakeven} không phụ
thuộc vào chúng; chỉ có biên độ tuyệt đối là phụ thuộc. Cần thay các ước tính bằng báo
giá trước bất kỳ quyết định mua sắm nào.
\end{remark}
